\scanpage{16}
\mspage{15}

Ceci dit, une figure $F'$ est une \uuline{subdivision} d'une figure $F$ ssi elle \uline{raffine} $F'$, et \inserted{si} \uline{elle a \unsure{$=$|$\overline{\mathrm{m}}$|mê} support}. \inserted{Mais (cf plus bas) cela équivaut : dire que $F$ est \deleted{\unsure{maximal|maximale}} \inserted{maximale} pour la relation $\leq$, \deleted{parmi} \deleted{\unsure{toutes les|tous les}} les figures qui raffinent $F$.}

Ainsi, on trouve d'autres éléments de structure :

\textcircled{5} Une relation d'ordre $\ll$ sur $\mathcal{F}_{el}$ (``raffinement''), impliquée par $\leq$, d'où une notion, pour deux figures $F\subset\mathcal{F}_{el}$ et $F'\subset\mathcal{F}_{el}$, de $F$ raffine $F'$ :

\[
\begin{aligned}
&\deleted{\unsure{\exists|F}}\ \forall\,\deleted{\unsure{F|X}}X\in\deleted{\unsure{\mathcal{F}_{el}|F_{el}}}F,\quad \exists X'\in F'\quad\text{telle que}\\
&X\ll X'.
\end{aligned}
\]

Ceci posé

\inserted{\uline{Fig}$_{\deleted{\unsure{2}}3}$ Pour $X,X'\in\mathcal{F}_{el}$, on : $X\leq X'$ ssi $\{X,X'\}$ comp. et $X\ll X'$.}

% slanted
\margin{Donc les structures \textcircled{3} et \textcircled{5} impliquent la structure \textcircled{1}}

\deleted{\textcircled{6} \unsure{1|I} : On appellera \unsure{de|la} subdivision d'une figure}

\inserted{\uline{Conséquence} Pour deux \unsure{\deleted{figures $F,F'$}|figures $F,F'$}$\subset\mathcal{F}_{el}$, on a $F\leq F'\Longleftrightarrow F,F'$ compatibles, et $F\ll F'$}

Ceci posé, on dira qu'une figure $F'$ est une \uline{subdivision} \unsure{de' une|d'une} autre $F$, si c'est un \uline{raffinement}, et si elle est maximale, pour la relation $\subset$, parmi les raffinements de $F$.
